First, relative wind components are computed in
aircraft frame :
\begin{eqnarray*}
U_a & = & - V_p.D^{-1} \\
V_a & = & - V_p.D^{-1}.\tan\beta + L.Vr_z\\
W_a & = & - V_p.D^{-1}.\tan\alpha - L.Vr_y
\end{eqnarray*}
where~: $ D = \sqrt{1+\tan^2\alpha+\tan^2\beta}$

Then, relative wind components are rotated
to the local geographical frame :
\newcommand\relphantom[1]{\mathrel{\phantom{#1}}}
\begin{equation*}\begin{split}
U_{ag} &= U_a.\sin\psi.\cos\theta +  
          V_a.(\cos\psi.\cos\phi + \sin\psi.\sin\theta.\sin\phi) \\
      &\relphantom{=} +  W_a.(\sin\psi.\sin\theta.\cos\phi - \cos\psi.\sin\phi) \\   
V_{ag} &= U_a.\cos\psi.\cos\theta - 
          V_a.(\sin\psi.\cos\phi - \cos\psi.\sin\theta.\sin\phi) \\
      &\relphantom{=}  + W_a.(\cos\psi.\sin\theta.\cos\phi + \sin\psi.\sin\phi) \\
W_{ag} &= U_a.\sin\theta - 
          V_a.\cos\theta.\sin\phi -
          W_a.\cos\theta.\cos\phi
\end{split}\end{equation*}
Eventually, the wind vector is the sum of the relative wind
and the platform speed relative to ground :
\begin{equation*}
\begin{pmatrix}U \\ V \\ W  \end{pmatrix} = 
\begin{pmatrix}U_{ag} \\ V_{ag} \\  W_{ag} \end{pmatrix} + 
\begin{pmatrix} V_e \\ V_n \\ V_v  \end{pmatrix}
\end{equation*}
